\section{Existence and uniqueness for measurable daughter kernels}
\label{app:existence}

This appendix proves \cref{thm:existence}. In addition to providing a
solution class for the estimates in the paper, the argument explains why
mere measurability of the parent-dependent daughter law causes no gap.
All the approximations converge in weighted variation, which controls
bounded measurable tests directly.

\subsection{Weak kernel integrals and a comparison lemma}

Throughout this appendix, $w(x)=1+x$, so that
$\norm{\mu}_w=\int w\dd|\mu|$ is the weighted variation \eqref{eq:variation}.
If $G_t$ is a signed measure kernel with
$\int_0^T\norm{G_t}_w\dd t<\infty$, its integral is the measure defined
on Borel sets by integration in time. Countable additivity follows from
dominated convergence for the variation kernels, and
\begin{equation}
 \norm{\int_s^tG_u\dd u}_w\leq\int_s^t\norm{G_u}_w\dd u.
 \label{eq:kernel-integral}
\end{equation}
The same convention applies to other positive weights.
This construction does not require strong measurability of $G_t$ as a
Banach-space-valued function. For instance, $t\mapsto\delta_t$ need not
be strongly measurable in total variation, even though all its time
integrals on Borel sets are well defined. We therefore use weak kernel
integrals, rather than a Bochner differential equation on measures.

We first establish the norm estimate used below. Suppose a signed
measure curve $d_t$ satisfies, in weak kernel form,
\begin{equation}
 d_t=d_0+\int_0^t\{H_s-a_s d_s\}\dd s,
 \label{eq:signed-loss}
\end{equation}
where $a_s(x)\geq0$ is jointly measurable and, for every finite $R,T>0$,
there is a nonnegative $h_{R,T}\in L^1(0,T)$ such that
$\sup_{0<x\leq R}a_s(x)\leq h_{R,T}(s)$ for almost every $s\in(0,T)$.
Assume also that $t\mapsto d_t(A)$ is measurable for each Borel $A$ and
that both $H_s$ and $a_sd_s$ have integrable weighted variation on $(0,T)$.
The loss integrating factor gives
\begin{equation}
 d_t(\dd x)=e^{-\int_0^t a_u(x)\dd u}d_0(\dd x)
   +\int_0^t e^{-\int_s^t a_u(x)\dd u}H_s(\dd x)\dd s.
 \label{eq:loss-duhamel}
\end{equation}

Here is the proof. By Tonelli's theorem, $(s,x)\mapsto\int_s^ta_u(x)\dd u$
is jointly measurable, and it is finite for every $x$ because
$a_\cdot(x)\leq h_{x,T}$ on $(0,T)$. The exponential factors are therefore
Borel functions with values in $[0,1]$, and the right-hand side of
\eqref{eq:loss-duhamel} is a finite signed measure on $E$.

Variations of setwise measurable curves are measurable in time, by a
countable supremum: if $t\mapsto\rho_t(A)$ is measurable for each Borel
$A$, then for every Borel $A$ and bounded Borel $\phi\geq0$,
\begin{equation}
 \int_A\phi\dd|\rho_t|
 =\sup_{B\in\mathcal B_0}\Bigl[\int_{A\cap B}\phi\dd\rho_t
                          -\int_{A\setminus B}\phi\dd\rho_t\Bigr],
 \label{eq:variation-countable-sup}
\end{equation}
where $\mathcal B_0$ is the countable algebra of finite unions of intervals
with rational endpoints. Indeed each bracket is at most the left side,
and the Hahn sets of $\phi\rho_t$ are approximated in
$\phi|\rho_t|$-measure by sets from the generating algebra $\mathcal B_0$.
Monotone convergence in $\phi$ extends this to unbounded Borel
$\phi\geq0$, in particular to $\norm{\rho_t}_w$.

Fix $R,T$ and restrict the output coordinate to $(0,R]$: for Borel
$A\subset(0,R]$, \eqref{eq:signed-loss} reads
$d_t(A)=d_0(A)+\int_0^t\{H_s(A)-(a_sd_s)(A)\}\dd s$. The restriction
applies only to the output coordinate; it does not discard any
contribution to $H_s$ from outside that set. Two facts follow from the
assumptions. First, the restricted curve is bounded in variation:
by \eqref{eq:kernel-integral} with weight one,
$\sup_{t\leq T}|d_t|((0,R])\leq\norm{d_0}_w
+\int_0^T(\norm{H_s}_w+\norm{a_sd_s}_w)\dd s<\infty$.
Second, on $(0,R]$ the loss is a multiplication operator bounded by an
integrable function: $|a_s\rho|((0,R])\leq h_{R,T}(s)|\rho|((0,R])$ for
every finite signed measure $\rho$ on $(0,R]$ and almost every $s$.
Consider the linear Volterra equation
\begin{equation}
 \rho_t(A)=d_0(A)+\int_0^tH_s(A)\dd s-\int_0^t(a_s\rho_s)(A)\dd s,
 \qquad A\subset(0,R]\ \text{Borel},\quad t\leq T,
 \label{eq:restricted-volterra}
\end{equation}
in the class of curves $\rho$ of finite signed measures on $(0,R]$ such
that $t\mapsto\rho_t(A)$ is measurable for each Borel $A$ and
$\sup_{t\leq T}|\rho_t|((0,R])<\infty$. The restriction of $d$ belongs to
this class and solves \eqref{eq:restricted-volterra}. Uniqueness holds in
this class: the difference $\delta$ of two solutions satisfies
$\delta_t=-\int_0^ta_s\delta_s\dd s$, hence
$|\delta_t|((0,R])\leq\int_0^th_{R,T}(s)|\delta_s|((0,R])\dd s$, and the
bounded function $s\mapsto|\delta_s|((0,R])$, measurable by
\eqref{eq:variation-countable-sup}, vanishes identically by Gronwall's
inequality. For existence, let $\rho_t$ be the restriction to $(0,R]$ of
the right-hand side of \eqref{eq:loss-duhamel}. It is a finite signed
measure with $|\rho_t|((0,R])\leq|d_0|(E)+\int_0^T|H_s|(E)\dd s$, and
$t\mapsto\rho_t(A)$ is measurable as an integral of a jointly measurable
integrand. For Borel $A\subset(0,R]$, the scalar identity
\[
 \int_r^ta_s(x)e^{-\int_r^sa_u(x)\dd u}\dd s=1-e^{-\int_r^ta_u(x)\dd u},
\]
valid because $a_\cdot(x)$ is integrable on $(0,T)$, and Fubini's theorem,
applicable because the integrands are dominated by $h_{R,T}(s)|d_0|(A)$
and $h_{R,T}(s)|H_r|(A)$, give
\begin{align*}
 \int_0^t(a_s\rho_s)(A)\dd s
 &=\int_A\bigl[1-e^{-\int_0^ta_u(x)\dd u}\bigr]d_0(\dd x)
  +\int_0^t\int_A\bigl[1-e^{-\int_r^ta_u(x)\dd u}\bigr]H_r(\dd x)\dd r\\
 &=d_0(A)+\int_0^tH_r(A)\dd r-\rho_t(A).
\end{align*}
So $\rho$ solves \eqref{eq:restricted-volterra}, and by uniqueness
$\rho_t$ is the restriction of $d_t$ to $(0,R]$. Since both sides of
\eqref{eq:loss-duhamel} are finite signed measures on $E$, letting
$R\to\infty$ through the integers and using countable additivity gives
the identity on every Borel set. This proves \eqref{eq:loss-duhamel}
for every curve $d$ satisfying the assumptions of this paragraph; in
particular, uniqueness is asserted among setwise measurable curves of
finite signed measures whose restrictions to bounded sets are bounded in
variation on bounded time intervals.

Define the nonnegative majorant
\begin{equation}
 z_t(\dd x)=e^{-\int_0^t a_u(x)\dd u}|d_0|(\dd x)
   +\int_0^t e^{-\int_s^t a_u(x)\dd u}|H_s|(\dd x)\dd s.
 \label{eq:z-majorant}
\end{equation}
Then $|d_t|\leq z_t$. Tonelli's theorem and the scalar loss identity
give
\begin{equation}
 \pair{w}{z_t}+\int_0^t\pair{wa_s}{z_s}\dd s
 =\norm{d_0}_w+\int_0^t\norm{H_s}_w\dd s.
 \label{eq:z-balance}
\end{equation}
In particular the integrated loss of $z$ is finite, without any higher
moment assumption on $z$.

\begin{lemma}[Weighted comparison]\label{lem:weighted-comparison}
In the preceding setting, suppose $P_s$ is a positive linear kernel
operator and $r_s$ a signed kernel such that
\begin{equation}
 |H_s|\leq P_s|d_s|+|r_s|,\qquad
 P_s^*w(x)-a_s(x)w(x)\leq c(s)w(x),
 \label{eq:comparison-assumptions}
\end{equation}
where $c\geq0$ is locally integrable, $s\mapsto\pair{P_s^*w-a_sw}{|d_s|}$
is measurable, and all terms paired with $|d_s|$ have finite time
integrals. Then
\begin{equation}
 \norm{d_t}_w\leq e^{\int_0^t c(u)\dd u}
 \left[\norm{d_0}_w+\int_0^t\norm{r_s}_w\dd s\right].
 \label{eq:comparison}
\end{equation}
\end{lemma}

\begin{proof}
Since $z_s\geq|d_s|$, \eqref{eq:z-balance} implies
\begin{align*}
 \norm{d_t}_w
 &\leq\norm{d_0}_w+
 \int_0^t\pair{P_s^*w-a_sw}{|d_s|}\dd s
 +\int_0^t\norm{r_s}_w\dd s\\
 &\leq\norm{d_0}_w+
 \int_0^t c(s)\norm{d_s}_w\dd s+
 \int_0^t\norm{r_s}_w\dd s.
\end{align*}
The function $s\mapsto\norm{d_s}_w$ is measurable by
\eqref{eq:variation-countable-sup} and locally bounded by
\eqref{eq:z-balance}, so Gronwall's inequality yields
\eqref{eq:comparison}.
Thus no differentiation of a norm or of a Hahn sign is needed.
\end{proof}

\subsection{Stability for the additive equation}

Let $n,v$ be two nonnegative curves with locally bounded count and second
moment, satisfying the population equation with the same coefficients,
possibly with an additional signed forcing term $r_t$ in their difference.
Put $d=n-v$ and $\mu=(n+v)/2$. Denote the moments of $\mu$ by
$\bar N,\bar m,\bar M_2$; the two curves need not have equal mass.
Polarization of the quadratic coagulation operator gives, against a test
$f$,
\begin{equation}
 \pair{f}{Q_t(n)-Q_t(v)}
 =\lambda(t)\iint (x+y)\Delta f(x,y)\mu_t(\dd x)d_t(\dd y).
 \label{eq:polarization}
\end{equation}
Write this difference and the fragmentation difference as
$H_t-a_td_t$, where
\begin{equation}
 a_t(y)=\lambda(t)(\bar m(t)+\bar N(t)y)+\sigma(t).
 \label{eq:comparison-loss}
\end{equation}
The terms in $H_t$ are the coagulation gain, the signed cross loss
$-\lambda(t)\mu_t(\dd x)\int(x+y)d_t(\dd y)$, the fragmentation gain,
and $r_t$. Their variations are bounded by $P_t|d_t|+|r_t|$, where
for a nonnegative measure $\xi$ and Borel $A$,
\begin{align}
 (P_t\xi)(A)
 &=\lambda(t)\iint(x+y)
       [\1_A(x+y)+\1_A(x)]\mu_t(\dd x)\xi(\dd y)\notag\\
 &\quad+\sigma(t)\int b_{t,y}(A)\xi(\dd y).
 \label{eq:positive-comparison}
\end{align}
Although the cross loss is replaced by a positive term, keeping the
direct loss $a_td_t$ preserves the cancellation needed for stability:
\begin{align}
 P_t^*w(y)-a_t(y)w(y)
 &=\lambda(t)\int(x+y)[w(x+y)+w(x)-w(y)]\mu_t(\dd x)
       +\sigma(t)\notag\\
 &=\lambda(t)[\bar m+2\bar M_2+y(\bar N+2\bar m)]+\sigma(t)
 \leq c(t)w(y),\label{eq:weight-cancellation}
\end{align}
where we may choose
\begin{equation}
 c(t)=\lambda(t)[\bar N(t)+3\bar m(t)+2\bar M_2(t)]+\sigma(t).
 \label{eq:stability-coefficient}
\end{equation}
All event integrals in weighted variation are finite under the stated
second-moment bounds. For example, the absolute weighted coagulation
increment is bounded by
$(x+y)[w(x+y)+w(x)+w(y)]=(x+y)[3+2(x+y)]$.
Also $a_t(y)\leq\lambda(t)(\bar m(t)+R\bar N(t))+\sigma(t)$ for
$0<y\leq R$, providing the required integrable local bound. The
measurability hypothesis of \cref{lem:weighted-comparison} holds because
$P_t^*w-a_tw$ in \eqref{eq:weight-cancellation} is a polynomial in $y$
whose coefficients are $\lambda(t)$, $\sigma(t)$, and moments of the
setwise measurable curve $\mu_t$; more generally, for a product-type
kernel $g(x,y)=\sum_i\phi_i(x)\psi_i(y)$ with finitely many Borel terms,
$s\mapsto\iint g\,\mu_s(\dd x)\,d_s(\dd y)$ and the same integral against
$|d_s|$ are finite sums of products of measurable functions of $s$, by
setwise measurability and \eqref{eq:variation-countable-sup}. All kernels
in \eqref{eq:polarization} and \eqref{eq:positive-comparison} that are
paired with $w$ are of this type.
\Cref{lem:weighted-comparison} therefore proves the stability estimate
\begin{equation}
 \norm{n_t-v_t}_w\leq
 e^{\int_0^t c(s)\dd s}
 \left[\norm{n_0-v_0}_w+\int_0^t\norm{r_s}_w\dd s\right].
 \label{eq:nonlinear-stability}
\end{equation}

Solutions in the sense of \cref{def:solution} are defined through
bounded tests only. \Cref{lem:weighted-comparison} needs
\eqref{eq:signed-loss} only on Borel sets, which is \eqref{eq:weak} with
indicator tests applied to $n$ and $v$ and subtracted; the integrability
of the weighted variations of $H_s$ and $a_sd_s$ follows from the locally
bounded second moments and the absolute event bound just displayed. The
identities for tests with $|f|\leq w$, obtained from the Borel-set form by
truncating $f$ and dominated convergence with that bound, enter only
through \eqref{eq:kernel-integral}, which gives the weighted-variation
continuity of $t\mapsto n_t$. In particular
\eqref{eq:nonlinear-stability} proves uniqueness in the solution class of
\cref{thm:existence} once existence is established.

\subsection{Bounded kernels with a finite third initial moment}

First assume $\int(1+x^3)n_0(\dd x)<\infty$. For integers $r\geq1$, replace
$x+y$ in the coagulation rate by
\begin{equation}
 k_r(x,y)=\min(x+y,r),
 \label{eq:cutoff}
\end{equation}
while retaining offspring size $x+y$ and the original fragmentation
kernel. In the space with norm $\norm{\mu}_3=\int(1+x^3)\dd|\mu|$, the
resulting coagulation vector field is locally Lipschitz, with a time
coefficient bounded on each ball by a constant times $r\lambda(t)$.
Indeed $1+(x+y)^3\leq4[(1+x^3)+(1+y^3)]$, and the bilinear difference
identity gives the estimate. Fragmentation is a bounded operator with
norm at most $3\sigma(t)$, since
\begin{equation}
 \int_E(1+z^3)b_{t,x}(\dd z)\leq2+x^3\leq2(1+x^3).
 \label{eq:fragment-weight}
\end{equation}

Picard iteration using weak kernel integrals constructs a local solution
continuous in this norm, on a time interval whose length depends only on
$r$, the rates, and the radius $H$ of the ball in $\norm{\cdot}_3$
containing the initial measure. Here are the relevant measurability and
positivity details. A continuous weighted-variation path is a measurable measure
kernel; composing it with the given daughter kernel and the coagulation
map preserves setwise measurability. Inequality \eqref{eq:kernel-integral}
makes each integral iterate continuous in norm and gives the usual
contraction estimate with an integrable time coefficient. On a nonnegative
ball of radius $H$, the loss rate is at most
$r\lambda(t)H+\sigma(t)$. Add this scalar multiple of the measure to the
vector field and solve its subtraction by an integrating factor. The
remaining map has nonnegative gain, maps a sufficiently short-time
nonnegative ball into itself, and is a contraction there. Since the
existence time depends only on $H$, $r$, and the rates, the construction
can be restarted from any time at which $\norm{n_t^r}_3$ is finite.
Iterating it gives a nonnegative maximal solution $n^r$, whose existence
time is either infinite or a time at which $\norm{n_t^r}_3$ is unbounded
on the approach.

The tests $1,x,x^2,x^3$ are valid for this bounded kernel in the weighted
space just used. Conservation of daughter mass and exact coagulation
mass cancellation yield
\begin{equation}
 M_1^r(t)=m,\qquad N^r(t)\leq N_0e^{F(t)}.
 \label{eq:cutoff-countmass}
\end{equation}
Since $z<x$ and the daughter first moment is $x$,
$\int z^q b_{t,x}(\dd z)\leq x^q$ for $q\geq1$.
Using $k_r\leq x+y$, the polynomial coagulation increments give
\begin{align}
 (M_2^r)'&\leq2\lambda(t)mM_2^r,\label{eq:cutoff-M2}\\
 (M_3^r)'&\leq3\lambda(t)\bigl[mM_3^r+(M_2^r)^2\bigr]
            \leq6\lambda(t)mM_3^r.\label{eq:cutoff-M3}
\end{align}
The final inequality follows from $(M_2^r)^2\leq mM_3^r$ by
Cauchy--Schwarz. Thus on each finite interval $[0,T]$ the count, second
moment, and third moment are bounded uniformly in $r$. Since
$\norm{n_t^r}_3=N^r(t)+M_3^r(t)$ then stays bounded on every finite
interval, the local existence time is bounded below along the maximal
solution, which therefore extends to the whole half-line.

Let $Q$ denote the full additive coagulation operator and $Q^r$ its
cutoff. Treat $n^r$ as a solution of the full equation with forcing
$e^r=Q^r(n^r)-Q(n^r)$. With $s=x+y$ its weighted norm satisfies
\begin{align}
 \norm{e_t^r}_w
 &\leq\frac{\lambda(t)}2\iint s\1_{\{s>r\}}(3+2s)
                       n_t^r(\dd x)n_t^r(\dd y)\notag\\
 &\leq\frac{\lambda(t)}{2r}\iint(3s^2+2s^3)
                       n_t^r(\dd x)n_t^r(\dd y).
 \label{eq:cutoff-error}
\end{align}
The two double moments are
$2N^rM_2^r+2m^2$ and $2N^rM_3^r+6mM_2^r$ respectively.
They are uniformly bounded on $[0,T]$. Applying
\eqref{eq:nonlinear-stability} to $n^r,n^q$, with forcing $e^r-e^q$,
therefore gives
\begin{equation}
 \sup_{t\leq T}\norm{n_t^r-n_t^q}_w
 \leq C_T(r^{-1}+q^{-1}),
 \label{eq:cutoff-Cauchy}
\end{equation}
where $C_T$ is independent of $r,q$. The coefficient in the stability
estimate is uniform because it requires only the common second-moment
bound, not a bound on the cutoff size.

Completeness gives a nonnegative limit $n_t$ in weighted variation,
uniformly on each finite interval. Its count and mass pass to the limit
in this norm; its second and third moments obey the preceding upper
bounds by monotone truncation and lower semicontinuity. For every bounded
Borel test, the coagulation map is continuous in this convergence:
polarization and $|\Delta f|\leq3\norm{f}_\infty$ bound its difference
by a constant times $\lambda(t)\norm{n^r_t-n_t}_w$ when the background
counts and masses are bounded. Fragmentation is likewise continuous
because $|\mathcal F_t f|\leq3\norm{f}_\infty$.
Equation \eqref{eq:cutoff-error} removes the forcing. Dominated convergence
in time passes the complete weak equation, including for merely Borel
parent-dependent daughters. The resulting $n$ has all the properties
claimed in \cref{thm:existence} for this initial class.

\subsection{Finite second initial moment}

For the general initial condition of \cref{thm:existence}, let
$n_0^j=\1_{(0,j]}n_0$. Each has finite third moment, and
$\norm{n_0^j-n_0}_w\to0$. The solutions just constructed have masses
$m_j\leq m$, initial counts at most $N_0$, and the common bounds
\begin{equation}
 N^j(t)\leq N_0e^{F(t)},\qquad
 M_2^j(t)\leq M_2(0)e^{2m\int_0^t\lambda(s)\dd s}.
 \label{eq:initial-approx-bounds}
\end{equation}
The unforced estimate \eqref{eq:nonlinear-stability} now makes $n^j$
Cauchy uniformly on $[0,T]$ in weighted variation, with no uniform third
moment requirement. The same passage to bounded Borel tests gives a
solution with initial measure $n_0$, mass $m_j\to m$, and the
second-moment bound \eqref{eq:M2bound}. Uniqueness has already been proved.

No tightness argument is needed: the limits were taken in the Banach
space of finite signed measures on $E$ with the norm $\norm{\cdot}_w$, so
$n_t$ is itself a finite measure on $E$ with the limiting mass, and no
boundary atom can arise. Finally testing with $1$ gives \eqref{eq:count},
completing the proof of \cref{thm:existence}.
